\section{Preliminaries} \label{sec:sr}

We begin by recalling definitions and background results that will be used throughout, following \cite[Chapter~III]{Stanley96} and \cite{Athanasiadis16}. We work over a field $k$. In particular, all rings are commutative $k$-algebras and singular homology is computed with coefficients in $k$.    

\subsection{Triangulations of simplices} \label{sec:triangulations}
In this section only, for the purposes of providing context, we allow that the field $k$ may be finite, and the triangulation $\sigma \colon \Gamma \to 2^V$ is not necessarily quasi-geometric.

 

We recall the notion of a homology triangulation, following \cite{Athanasiadis12}.  A $d$-dimensional simplicial complex $\Gamma$ with trivial reduced homology is a \emph{homology ball} 
of dimension $d$ if there is a subcomplex $\partial \Gamma \subset \Gamma$  such that 
\begin{itemize}
\item $\partial \Gamma$ is a homology sphere of dimension $d -1$,
\item $\lk_\Gamma(F)$ is a homology sphere of dimension $d - |F|$ for $F \not \in \partial \Gamma$.
\item $\lk_\Gamma(F)$ is a homology ball of dimension $d - |F|$ for all nonempty $F \in \partial \Gamma$.
\end{itemize}
The \emph{interior faces} of a homology ball $\Gamma$ are the faces not contained in $\partial \Gamma$.  A  \emph{homology triangulation} of the simplex $2^V$ is a finite simplicial complex $\Gamma$ and a map $\sigma\colon \Gamma \to 2^V$ such that for every non-empty $U \subset V$,
\begin{itemize}
\item the simplicial complex $\Gamma_U  \coloneqq  \sigma^{-1}(2^U)$ is a homology ball of dimension $\vert U \vert - 1$.
\item $\sigma^{-1}(U)$ is the set of interior faces of the homology ball $\sigma^{-1}(2^U)$.
\end{itemize}
Note that the Betti numbers of a simplicial complex, and hence the property of being a homology ball, depend only on the characteristic of the field $k$. Homology triangulations are a special case of the (strong) formal subdivisions of Eulerian posets considered in \cite[Section~7]{Stanley92} and \cite[Section~3]{KatzStapledon16}.

The \emph{carrier} of a face $F \in \Gamma$ is $\sigma(F)$. A homology triangulation $\sigma\colon \Gamma \to 2^V$ is \emph{quasi-geometric} if there is no face $F \in \Gamma$ and $U \subset V$ such that the dimension of $\Gamma_U$ is strictly smaller than the dimension of $F$ and the carrier of every vertex in $F$ is contained in $U$.  A homology triangulation is \emph{geometric} if it can be realized in $\R^n$ as the subdivision of a geometric simplex into geometric simplices. Every geometric homology triangulation is quasi-geometric. 

The local $h$-vector, which we have defined in the introduction as the Hilbert function of the local face module, can be expressed in terms of $h$-vectors of subcomplexes of links of faces in the homology balls $\Gamma_U$:
\begin{equation} \label{eq:localh}
\ell(\Gamma, E) = \sum_{U \supset \sigma(E)} (-1)^{|V| - \vert U \vert} h(\lk_{\Gamma_U}(E)).
\end{equation}
Note that \eqref{eq:localh} makes sense even when $k$ is finite or $\sigma \colon \Gamma \to 2^V$ is not quasi-geometric, and should be taken as the definition of the local $h$-vector in this broader context. 


\begin{theorem}[\cite{Stanley92, Athanasiadis12, KatzStapledon16}]\label{t:localproperties}
Let $\sigma \colon \Gamma \to 2^V$ be a homology triangulation, let $E$ be a face of $\Gamma$ and let $d = |V| - |E|$.  Then the local $h$-vector $(\ell_0, \ldots, \ell_d)$ satisfies: \\
\begin{tabular}{lll}
\quad \quad $\bullet$ &  \emph{(symmetry)} &
$\ell_i = \ell_{d-i};$ \\
\quad \quad $\bullet$ & \emph{(non-negativity)} & if $\Gamma$ is quasi-geometric, then $\ell_i \geq 0;$\\
\quad \quad $\bullet$ & \emph{(unimodality)} &
 if $\Gamma$ is regular, then $\ell_0 \leq \ell_1 \leq \cdots \leq \ell_{\lfloor d/2 \rfloor}$.
\end{tabular}
\end{theorem}

Note that the proof of non-negativity for quasi-geometric triangulations, due to Stanley and Athanasiadis, is via the identification with the Hilbert function of the local face module. It suffices to consider the case where $k$ is infinite, since \eqref{eq:localh} is invariant under field extensions.


\subsection{Face rings and special l.s.o.p.s}\label{ss:sr}

Here, and for the remainder of the paper, the field $k$ is fixed and infinite, and all triangulations are quasi-geometric homology triangulations.

 Given a finite simplicial complex $\Gamma$ with vertex set $V = \{v_1, \ldots, v_n\}$, let $k[\Gamma]$ denote the \emph{face ring}. In other words, for each subset $F \subset V$, let $x^F$ be the corresponding squarefree monomial in the polynomial ring $k[x_1, \ldots, x_n]$, i.e.
 $
 x^F \coloneqq  \prod_{v_i \in F} x_i.
 $ 
 Then the face ring is
 \[
 k[\Gamma]  \coloneqq  k[x_1, \ldots, x_n] / (x^F : F \mbox{ is not a face in } \Gamma).
 \]
Given a subcomplex $\Gamma'$ of $\Gamma$, we have a natural restriction map $k[\Gamma] \rightarrow  k[\Gamma']$, taking $x^F$ to $x^F$ if $F \in \Gamma'$ and to 0 otherwise. Given $\theta \in k[\Gamma]$, let $\theta|_{\Gamma'}$ denote the image of $\theta$ in $k[\Gamma']$. In particular, each $F$ in $\Gamma$ may be viewed as a subcomplex, and we write $\theta|_F$ for the restriction of $\theta$ to this subcomplex.

Note that $k[\Gamma]$ is graded by degree. By definition, a linear system of parameters (l.s.o.p.) for a finitely generated graded $k$-algebra $R$ of Krull dimension $d$ is a sequence of elements $\theta_1, \ldots, \theta_d$ in $R_1$ such that $R/(\theta_1, \ldots, \theta_d)$ is a finite-dimensional $k$-vector space. If $\Gamma$ is a Cohen-Macaulay complex (i.e. if $k[\Gamma]$ is a Cohen-Macaulay ring) and $\theta_1, \ldots, \theta_d$ is an l.s.o.p.\ for $k[\Gamma]$, then $(\theta_1, \ldots, \theta_d)$ is a regular sequence and the $h$-polynomial of $\Gamma$ is the Hilbert series of $k[\Gamma]/(\theta_1, \ldots, \theta_d)$.  Links of faces in triangulations of simplices are Cohen-Macaulay \cite{Reisner76}. 

Suppose $\Gamma$ has dimension $d-1$, so $k[\Gamma]$ has Krull dimension $d$.  Then a sequence of elements $\theta_1, \ldots, \theta_d$ in $k[\Gamma]_1$ is an l.s.o.p. for $k[\Gamma]$ if and only if the following condition is satisfied \cite[Lemma~2.4(a)]{Stanley96}:


%\renewcommand{\labelitemi}{$(*)$}
\begin{enumerate}[label=(*)]
\item \label{it1} For every face $F \in \Gamma$ (or equivalently, for every facet $F \in \Gamma$), the restrictions $\theta_1|_F, \ldots, \theta_d|_F$ span a vector space of dimension $|F|$. 
\end{enumerate}

This characterization provides flexibility in constructing l.s.o.p.s in which the linear functions have specified support, where the \emph{support} of $\theta = \sum a_i x_i$ is $\supp(\theta)  \coloneqq  \{ v_i : a_i \neq 0 \}$.

\begin{lemma} \label{lemma:lsopexistence}
Let $S_1, \ldots, S_d$ be subsets of the vertices of $\Gamma$.  Then there is an l.s.o.p. $\theta_1, \ldots, \theta_d$ for $k[\Gamma]$ such that $\supp (\theta_i) = S_i$ for $1 \leq i \leq d$ if and only if, for every face $F \in \Gamma$,
\begin{equation} \label{eq:marriageinequality}
| \{ S_i : S_i \cap F \neq \emptyset \}| \geq |F|.
\end{equation}
\end{lemma}

\begin{proof} 
The argument is similar to that given by Stanley in \cite[Corollary~4.4]{Stanley92}. The necessity of \eqref{eq:marriageinequality} follows immediately from~\ref{it1}. We now prove its sufficiency.  Suppose $S_1, \ldots, S_d$ are chosen such that \eqref{eq:marriageinequality} holds for every 
$F \in \Gamma$.    
Let $N = |S_1| + \cdots + |S_d|$, and consider the space $k^N$ 
parametrizing tuples $(\theta_1, \ldots, \theta_d)$
with $\supp (\theta_i) \subset S_i$.
Fix $F = \{v_1, \dotsc, v_k\} \in \Gamma$.
Let $X_F \subset k^N$  parametrize the tuples 
 whose restrictions to $F$ 
 span a vector space of dimension $|F|$.
Note that $X_F$ is Zariski open.  
By Hall's Marriage Theorem, there is a permutation $\sigma \in \mathfrak{S}_d$ such that $v_i \in S_{\sigma(i)}$. If we set $\theta_{\sigma(i)} = x_i$ for $1 \le i \le k$, and $\theta_{\sigma(i)} = 0$ for $i > k$, then $\theta \in X_F$, and hence $X_F$ is nonempty. Also, the subset of $k^N$ where all coordinates are nonzero is Zariski open and nonempty.  Since $k$ is infinite, the intersection of these nonempty Zariski open subsets of $k^N$ is nonempty, and hence there is an l.s.o.p. $\theta_1, \ldots, \theta_d$ with $\supp(\theta_i)= S_i$. 
\end{proof}




Let $\sigma \colon \Gamma \to 2^V$ be a quasi-geometric homology triangulation,   and let $E \in \Gamma$ be a~face. 

\begin{definition}[\cite{Stanley92, Athanasiadis12b}] \label{def:special} A linear system of parameters $\theta_1, \dotsc, \theta_{d}$ for $k[\lk_{\Gamma}(E)]$ is \textit{special} if, for each vertex $v \in V$ with $v \not \in \sigma(E)$, there is an element $\theta_v$ of the l.s.o.p. such that $\supp(\theta_v)$ consists of vertices in $\lk_{\Gamma}(E)$  whose carrier contains $v$, and such that $\theta_v \not= \theta_{v'}$ for $v \not= v'$. 
\end{definition}

In other words, after reordering so that $\sigma(E)^c = \{v_1, \ldots, v_b\}$, an l.s.o.p. for $k[\lk_\Gamma(E)]$ is special if we can order it $\theta_1, \ldots, \theta_d$ such that
\[
\supp(\theta_i) \subset \{ w \in \lk_\Gamma(E) : v_i \in \sigma(w)\},
\]
for $1 \leq i \leq b$.  The existence of special l.s.o.p.s is well-known to experts and the proof is similar to Stanley's argument in the case $E = \emptyset$. For completeness, we provide a short proof.

\begin{proposition} \label{prop:speciallsopexistence}
Suppose $k$ is infinite. Let $\sigma\colon \Gamma \to 2^V$ be a quasi-geometric homology triangulation of a simplex, and let $E$ be a face of $\Gamma$. Then there is a special l.s.o.p. for $k[\lk_{\Gamma}(E)]$.
\end{proposition}
\begin{proof}

Let $V = \{v_1, \ldots, v_n\}$.  After renumbering, we may assume that $\sigma(E)^c = \{v_1, \dotsc, v_b\}$. Fix $d = n - \vert E \vert$. Note that $b \leq d$.  We define subsets $S_1, S_2, \dotsc, S_{d}$ of the vertices in $\lk_{\Gamma}(E)$, as follows. For $i \le b$, let $S_i$ be the set of vertices $w$ such that $v_i \in \sigma(w)$. 
For $i > b$, let $S_i$ be the set of all vertices of $\lk_{\Gamma}(E)$. Because $\sigma$ is quasi-geometric, for each face $F$ of $\lk_{\Gamma}(E)$, the union of the sets $\sigma(w) \subset V$, as $w$ ranges over vertices of $E \sqcup F$, has size at least $|E| + |F|$.  It follows that $|\{i \leq b : S_i \cap F \neq \emptyset \}| \geq |F| - (d-b)$.  Since $S_j \cap F \neq \emptyset$ for $j > b$, we conclude that $|\{i : S_i \cap F \neq \emptyset \}| \geq |F|$.  Hence, by Lemma~\ref{lemma:lsopexistence}, there is an l.s.o.p. $\theta_1, \ldots, \theta_d$ for $k[\lk_\Gamma(E)]$ with $\supp(\theta_i) = S_i$. 
\end{proof}
